\section{Discussion}
\label{sec:discussion}

The experiments in Section~\ref{sec:experiments} demonstrate that the
\textsc{ComposeHead} can recover exact closed-form solutions to the
two-dimensional Navier--Stokes equations from data.  We now examine what makes
the pipeline work, situate it relative to existing methods, and identify the
limitations that constrain its present applicability.

\paragraph{What the method can do.}
The method discovers closed-form PDE solutions whenever the true solution can be
expressed as a linear combination of products of $\sin$, $\cos$, $\exp$, $\ln$,
and the identity applied along individual coordinate axes.  This class
encompasses a substantial fraction of the analytical solutions in classical
references: separable solutions to heat, wave, and advection--diffusion
equations; Fourier modes of elliptic boundary-value problems; and exponentially
decaying vortex solutions of the Navier--Stokes equations.  The
ansatz~\eqref{eq:composehead} aligns naturally with the classical method of
separation of variables, a primary tool for constructing exact PDE solutions for
over three centuries.  The novelty is not the separable form itself but the
end-to-end pipeline from data to exact symbolic expressions without requiring
human insight into the solution structure.  The method also extends to
non-separable solutions that decompose as finite sums of separable terms.  The
traveling wave $\sin(x - 2t) = \sin(x)\cos(2t) - \cos(x)\sin(2t)$ was
recovered exactly via a brute-force two-term pair search, demonstrating that
the effective hypothesis class is strictly larger than the set of single
separable products.

\paragraph{What the method cannot do.}
The \textsc{ComposeHead} does not prove existence or regularity of solutions in
the PDE-theoretic sense, does not address uniqueness, and does not guarantee
that a discovered solution is the physically relevant one.  Turbulent flows lie
entirely outside the method's reach.  More generally, any solution not
expressible as a finite linear combination of separable products of the
available primitives is undiscoverable.

\paragraph{Relationship to the Navier--Stokes Millennium Prize.}
The recovery of the Taylor--Green vortex naturally raises the question of how
this work relates to the Clay Mathematics Institute Millennium Prize Problem.
We wish to be precise: the Prize requires proving, for the three-dimensional
incompressible Navier--Stokes equations, that smooth solutions exist globally in
time for arbitrary smooth initial data of finite energy, or exhibiting a
counterexample.  Our work addresses a fundamentally different question: given
specific initial conditions, can one algorithmically discover a closed-form
expression that satisfies the equations?  The result is a \emph{constructive
existence demonstration} for a particular initial datum---we exhibit the
solution rather than proving one must exist.  This is distinct from, and
considerably weaker than, the general existence claim demanded by the Prize.
Nonetheless, algorithmically discovering closed-form solutions for families of
initial conditions may inform theoretical investigations by providing exact
solutions whose analytic properties can be studied directly.

\paragraph{The role of normalization.}
The normalization pipeline is the central practical contribution of this work.
The parameterization $c_k \cdot f(\alpha_{k,d}\, x_d + \beta_{k,d})$ admits
gauge freedom: multiple parameter settings encode the same function.  For
example, $\sin(x - \pi) \cdot \exp(b) \cdot c = \sin(x) \cdot c'$ where
$c' = -\exp(b) \cdot c$.  A cosine with phase $\pi/2$ is a sine; exponential
biases can be absorbed into multiplicative coefficients.  Without
normalization, gradient-based optimization converges to one of many equivalent
settings, and coefficient snapping fails.  The normalization pipeline resolves
this gauge freedom by absorbing exponential biases into~$c_k$, canonicalizing
trigonometric phases via $\sin(x + \pi/2) = \cos(x)$ and
$\cos(x + \pi) = -\cos(x)$, and reducing all phase offsets to $[0, \pi/2)$.
This canonicalization transforms a continuously parameterized model into an
expression that coefficient snapping can convert to exact symbolic form,
reducing parameter errors from arbitrary gauge-equivalent configurations to the
order of $10^{-6}$.

\paragraph{Single-term search.}
Equally important is searching for one term at a time rather than fitting a
multi-term linear combination simultaneously.  Joint optimization allows the
optimizer to split a single physical signal across several basis
elements---representing $\sin(x)\cos(y)$ as a sum of two terms with
compensating coefficients.  Such splitting yields numerically accurate fits but
expressions bearing no resemblance to the true solution.  By fitting a single
term, subtracting it from the residual, and repeating, each term is forced to
capture a complete, interpretable component.  This greedy strategy is a
practical choice rather than a theoretical necessity, but it proves essential
for exact, human-readable recovery.

\paragraph{What EML provides (and what it does not).}
The Euler--Maclaurin--Lie framework contributes two things.  First, it provides
a principled, non-ad-hoc justification for the primitive basis
$\{\sin, \cos, \exp, \ln\}$: these arise as fundamental solutions of linear
constant-coefficient differential equations and their multiplicative
counterparts.  Without EML the same library could be chosen on empirical
grounds, but the basis selection would lack theoretical grounding.  Second, the
EML basis is closed under differentiation, guaranteeing that PDE residuals
computed from a candidate expression remain within the representable class and
that residual-based loss functions stay well-defined throughout optimization.

The normalization pipeline, however, does not follow from EML per se.  The
gauge-resolving identities exploit specific algebraic properties of $\exp$,
$\sin$, and $\cos$ that hold independently of the EML framework.  EML
determines \emph{which} functions to include; normalization determines
\emph{how} to make optimization over those functions yield exact symbolic
output.

EML universality motivated the architectural choices that created the
conditions under which normalization and snapping could succeed.  The
conceptual chain---universality $\to$ verified primitives $\to$ compositional
head $\to$ structured parameters $\to$ normalization $\to$ exact
recovery---starts with EML\@.  However, the normalization pipeline itself
exploits properties of $\exp$, $\sin$, and $\cos$ that are independent of
the EML implementation, and the numerical results would be identical if the
primitives were implemented via PyTorch intrinsics rather than EML
compositions.  EML universality motivates the framework and guarantees
completeness of the primitive basis, but the practical
contributions---normalization and search strategy---are independent of the EML
implementation and would apply equally to any library of elementary functions.

\paragraph{EML branch closure: predicting solution existence.}
The EML-derived primitives naturally partition into three \emph{branches}
based on their behavior under differentiation and multiplication:
\textbf{Branch~E} ($\exp$), which is closed under both operations;
\textbf{Branch~T} ($\sin$, $\cos$), which is closed under differentiation
but generates new frequencies under multiplication
($\sin a \cdot \sin b \to \cos(a\!-\!b) - \cos(a\!+\!b)$); and
\textbf{Branch~L} ($\ln$), which is not closed under differentiation
($\mathrm{d}/\mathrm{d}x\;\ln x = 1/x$).  A PDE's nonlinear operator~$N$
maps solutions from one branch combination to another, and the existence of
exact separable solutions depends on whether this mapping \emph{closes}:
\begin{itemize}
  \item \textbf{Direct closure.}  If $N$ maps a branch back to itself
    (as linear operators do for Branch~T), exact separable solutions exist
    generically.  The heat and wave equations fall into this category.
  \item \textbf{Gradient closure.}  If $N$ generates terms outside the
    branch but those terms are a pure gradient---absorbable into an
    auxiliary variable such as pressure---the equation is effectively
    closed.  The 2-D Navier--Stokes equations satisfy this condition:
    the quadratic term $(\mathbf{u}\!\cdot\!\nabla)\mathbf{u}$ produces
    $\sin(2x)$ from $\sin(x)$, but this equals
    $\partial_x[-\tfrac{1}{4}\cos(2x)]$ and is absorbed into~$\nabla p$.
  \item \textbf{No closure.}  When neither mechanism applies, the
    nonlinearity generates irreducible new modes, and no finite
    separable solution exists.
\end{itemize}
This framework yields verifiable predictions.  We tested it on eight PDEs
(Table~\ref{tab:branch}): the Burgers, KdV, and nonlinear Schr\"{o}dinger
equations all fail closure in Branch~T, which the PDE-residual solver
confirms by producing residuals orders of magnitude above machine precision.
The 3-D non-Beltrami Navier--Stokes equations fail gradient closure
because vortex stretching introduces an irreducible term
(Section~\ref{sec:3d_ns}).  All four predictions of success and all four
predictions of failure were confirmed computationally, suggesting that
branch closure analysis can serve as a \emph{pre-screening tool}
that identifies which PDEs are amenable to separable symbolic
recovery before any optimization is attempted.

\begin{table}[t]
  \centering
  \caption{EML branch closure predictions and computational verification.
    ``Closure'' indicates the mechanism by which the nonlinear term stays
    within the separable hypothesis class.  All predictions were confirmed.}
  \label{tab:branch}
  \smallskip
  \begin{tabular}{@{}llcll@{}}
    \toprule
    PDE & Nonlinearity & Branch & Closure & Outcome \\
    \midrule
    1-D Heat       & Linear    & T & Direct   & Exact recovery \\
    1-D Wave       & Linear    & T & Direct   & Exact recovery \\
    2-D Navier--Stokes & Quadratic & T & Gradient & Exact recovery \\
    3-D NS (Beltrami)  & Quadratic & T & Beltrami & Exact recovery \\
    \midrule
    3-D NS (non-Beltrami) & Quadratic & T & Fails & Correct failure \\
    1-D Burgers    & Quadratic & T,E & Fails  & Correct failure \\
    1-D KdV        & Quadratic & T & Fails    & Correct failure \\
    NLS standing wave & Cubic  & T & Fails    & Correct failure \\
    \bottomrule
  \end{tabular}
\end{table}

\paragraph{Comparison to symbolic regression.}
PySR~\cite{cranmer2023} and AI~Feynman~\cite{udrescu2020} search over arbitrary
expression trees and are strictly more general: they can discover algebraic
relationships of essentially any form, including non-separable expressions such
as $E = mc^{2}$.  Our approach trades this generality for structure.  The
trade-off is justified by a concrete gap.  General symbolic regression methods
struggle with multi-variable separable products common in PDE solutions:
recovering $\sin(x)\cos(y)\exp(-t/5)$ requires simultaneously discovering the
correct univariate functions, their arguments, and the multiplicative
coupling---a combinatorially explosive task.  Physics-informed neural networks
avoid this search but produce weight matrices, not symbolic formulas.  The
\textsc{ComposeHead} fills this gap by constraining the hypothesis class to
separable products, where free parameters scale linearly in the number of terms
rather than exponentially in expression depth.

\paragraph{Limitations.}
\begin{itemize}
  \item \emph{Data dependence.}  The method requires training data from
    simulation, experiment, or collocation-point sampling; it cannot discover
    solutions from the PDE alone.
  \item \emph{Separability assumption.}  The ansatz~\eqref{eq:composehead}
    cannot represent non-separable solutions such as self-similar forms
    $u = t^{-\alpha} g(x/t^{\beta})$ or traveling waves $u = f(x - ct)$.
  \item \emph{Combinatorial scaling.}  The candidate count scales as
    $|\mathcal{F}|^{K} \times D^{K}$, exponential in the product order~$K$.
  \item \emph{Dimensionality.}  Experiments are restricted to two spatial
    dimensions plus time; three-dimensional validation remains future work.
  \item \emph{Basis completeness.}  The primitive library must contain the
    functions in the true solution.  Bessel functions, elliptic integrals, and
    other special functions are undiscoverable unless added to the basis.
\end{itemize}


\section{Future Work}
\label{sec:future}

Several directions extend the framework presented here.

\paragraph{Three-dimensional Navier--Stokes solutions.}
Exact 3D Navier--Stokes solutions---Beltrami flows, the Arnold--Beltrami--%
Childress (ABC) flow---have velocity fields that are separable products of
trigonometric functions in $x$, $y$, $z$.  These lie within the
\textsc{ComposeHead} hypothesis class and would directly test scalability to
higher-dimensional domains.

\paragraph{Physics-informed discovery without training data.}
A natural next step is to train using only the PDE residual, boundary
conditions, and initial conditions---in the spirit of physics-informed neural
networks but with a symbolic ansatz.  The differentiation closure of the EML
basis makes this extension particularly natural, since PDE residuals remain
within the representable function class.

\paragraph{Non-separable and nested compositions.}
Two strategies may extend the method beyond separability: (i)~mixed terms
$f(ax + by)$ coupling coordinate variables, and (ii)~nested compositions
$f \circ g$ within each factor.  Both enlarge the hypothesis class at the cost
of combinatorial complexity.

\paragraph{Extending the primitive basis: hyperbolic functions.}
The primitive basis $\{\sin, \cos, \exp, \ln\}$ can be extended with hyperbolic
functions $\tanh$ and $\operatorname{sech}$, which are EML-constructible by
universality---$\tanh(x) = (\exp(x) - \exp(-x))/(\exp(x) + \exp(-x))$ is a
finite composition of $\exp$ and division---but impractically deep as raw EML
trees.  Extracting them as named primitives is the mathematical equivalent of
defining a subroutine: compressing a deep composition into a stable,
differentiable unit.  We verified that this extension recovers the Allen--Cahn
kink $u = \tanh(x/\sqrt{2})$ (pointwise error $< 10^{-8}$), the Burgers steady
shock $u = A\tanh(Ax/2\nu)$, and the KdV soliton profile $u = -2\operatorname{sech}^{2}(x)$---solutions
entirely outside the reach of the original basis.  This pattern generalizes:
Bessel functions, Airy functions, and elliptic functions are all
EML-constructible at sufficient depth and could be ``unfolded'' into named
primitives to unlock their respective PDE classes.

\paragraph{Function discovery via unconstrained EML.}
A deeper question is whether EML can discover function families that have not
yet been named.  Historically, new functions were discovered when PDEs demanded
them: Bessel functions (1817) from the cylindrical wave equation, Airy functions
(1838) from optics, Weierstrass $\wp$ (1860s) from elliptic PDEs.  In each
case, the equation forced the creation of a new primitive.  The raw
\texttt{EMLHead}---a binary tree of \texttt{eml}$(x,y) = \exp(x) - \ln(y)$
nodes with trainable weights---searches over \emph{all} compositions of $\exp$
and $\ln$ up to a given depth, unconstrained by named functions.  If trained on
a PDE residual, the resulting tree \emph{is} the solution, expressed as a nested
composition.  If that composition does not simplify to any known function, it
constitutes a genuine discovery: a function defined by its EML structure.
Preliminary experiments demonstrate feasibility at multiple levels.  First,
population-based optimization with evolutionary selection recovers $\tanh$ from
raw EML at MAE~$\sim 3\!\times\!10^{-4}$ (correlation $r = 1.0000$) and
even approximates non-elementary functions: $\operatorname{erf}(x)$---defined
by an integral with no finite exp/ln representation---is recovered at
MAE~$\sim 1.6\!\times\!10^{-4}$.  Second, PDE-residual-only training
discovers $\tanh(x/\sqrt{2})$ from the Allen--Cahn equation without
knowing $\tanh$ exists (MAE~$\sim 1.9\!\times\!10^{-3}$).  Third, targeting
ODEs with \emph{no known closed-form solution}---the Blasius boundary-layer
equation $f''' + ff'' = 0$ (unsolved symbolically since 1908) and the
Lane--Emden $n\!=\!3$ polytrope $y'' + (2/x)y' + y^{3} = 0$---yields EML
trees that reproduce the numerical solutions at MAE~$\sim 10^{-3}$ and
$7\!\times\!10^{-4}$ respectively.  These trees are, in a precise sense,
the first symbolic representations of these functions: finite compositions of
$\exp(w \cdot x + b) - \ln|w \cdot y + b|$ that encode functions which have
resisted closed-form expression for over a century.

Cross-tree motif analysis reveals that the Blasius and Lane--Emden EML trees
share internal sub-compositions: both independently construct
$\ln(1 + \exp(x))$ (the softplus function) as a recurring building block,
with cross-tree correlation $r > 0.999$.  This suggests that EML trees
converge on common structural motifs even when fitting unrelated functions,
and that these motifs may represent a natural ``vocabulary'' of EML
sub-expressions.  However, the optimization landscape for deep compositions
remains orders of magnitude harder than for named primitives: deeper trees
(depth 5--6) with 188--380 parameters consistently underperform shallower
trees (depth 4, 92 parameters) when optimization budgets are fixed.
This confirms that naming a function is not merely a notational convenience
but a critical compression that makes optimization tractable.  Developing
better optimization methods for deep EML trees---CMA-ES with larger
populations, hierarchical motif extraction with progressive freezing, or
reinforcement-learning-guided tree construction---is a central open problem
that could transform EML from a tool for recovering known solutions into
a tool for discovering new mathematics.

\paragraph{Automated basis selection.}
An adaptive approach that begins with a small library and introduces additional
primitives based on residual analysis could reduce the combinatorial burden when
the relevant function class is unknown a priori.

\paragraph{Multi-term scaling.}
The brute-force pair search scales as $|\mathcal{F}|^{2K}$ for two-term sums,
which is tractable for small $K$ and $|\mathcal{F}|$ but becomes prohibitive
for higher-order decompositions.  Principled strategies for multi-term
discovery---greedy sequential fitting with gauge-aware normalization, or
structured sparsity over a large two-term dictionary---remain open problems.

\paragraph{Scaling to higher-dimensional PDEs.}
PDEs in kinetic theory, finance, and quantum mechanics may involve five or more
variables.  Algorithmic innovations---sparsity-promoting regularization,
hierarchical decomposition, or greedy term selection---will be needed to manage
the exponential growth of the candidate term count.

\paragraph{Integration with existing PDE solver libraries.}
The \textsc{ComposeHead} could serve as a post-processing module for numerical
solvers, accepting simulation data and returning symbolic expressions.  Hybrid
workflows alternating between numerical solution and symbolic identification may
prove effective for parameter studies.


\section{Conclusion}
\label{sec:conclusion}

We have presented the \textsc{ComposeHead}, a pipeline for discovering
closed-form solutions to partial differential equations from data.  The two
core practical contributions are: (i)~a normalization pipeline that resolves
the gauge freedom inherent in parameterizations of elementary functions,
collapsing families of equivalent parameter settings to canonical forms amenable
to coefficient snapping; and (ii)~a single-term greedy search strategy that
prevents signal splitting and ensures each discovered term corresponds to a
complete, interpretable component of the solution.  These components enable the
transition from a continuously parameterized model to an exact symbolic
expression.  The EML universal basis provides the theoretical foundation,
supplying a principled justification for the primitive function library and
guaranteeing differentiation closure so that PDE residuals remain within the
representable function class.

We demonstrated the method on the two-dimensional incompressible Navier--Stokes
equations, recovering the Taylor--Green vortex solution
$u = \sin(x)\cos(y)\exp(-t/5)$,\;$v = -\cos(x)\sin(y)\exp(-t/5)$
at machine precision with pointwise errors below $10^{-6}$ after normalization
and coefficient snapping.  The resulting expressions are genuine symbolic
formulas that can be analytically differentiated, substituted into the governing
equations, and verified to satisfy the Navier--Stokes system exactly.

The method fills a genuine gap.  General symbolic regression
methods~\cite{cranmer2023,udrescu2020} struggle with multi-variable separable
products; physics-informed neural networks produce weight matrices rather than
formulas.  By pairing a structured separable ansatz with gauge-resolving
normalization, the \textsc{ComposeHead} provides an automatic path from data to
exact symbolic PDE solutions---a capability that, to our knowledge, no existing
tool offers.

We release \texttt{torch-eml} as an open-source PyTorch library at
\url{https://github.com/gaius050526/torch-eml}, providing verified EML
primitive implementations, the \textsc{ComposeHead} architecture, and the
normalization and snapping pipeline described in this work.

While the current method is limited to separable solutions in two spatial
dimensions and requires the primitive basis to contain the relevant functions,
it establishes a foundation for symbolic PDE solving that occupies a distinctive
position between neural solvers and classical analytical methods.  Neural
solvers offer generality at the cost of interpretability; analytical methods
offer exactness but require human ingenuity.  The approach presented here
automates the discovery of exact, interpretable solutions within a structured
function class, bridging the gap between these two paradigms.
